% TOP_PROBLEM: {"id": 30006813, "title": "Dynamical Mordell–Lang Conjecture", "category_id": 1, "category": {"id": 1, "name": "number_theory", "display_name": "Number Theory", "description": "Properties of integers, prime numbers, Diophantine equations.", "slug": "number-theory", "order_index": 1, "created_at": "2026-07-31T15:26:25.670Z"}, "status": "open"}
% FIRST_POSTED: 2026-09-25
% ENTRY_KIND: statement_only
% !TeX program = lualatex
% Definition and sources only; this file is not an LLM solution attempt.
\documentclass[11pt]{article}
\usepackage[margin=25mm]{geometry}
\usepackage{fontspec}
\setmainfont{Latin Modern Roman}
\usepackage{amsmath,amssymb,mathtools}
\usepackage{unicode-math}
\setmathfont{Latin Modern Math}
\usepackage{xurl,hyperref}
\hypersetup{colorlinks=true,urlcolor=blue}
\setcounter{secnumdepth}{0}
\setlength{\parindent}{0pt}
\setlength{\parskip}{5pt}
\setlength{\emergencystretch}{3em}
\begin{document}
\section*{\texorpdfstring{Dynamical Mordell–Lang Conjecture}{Dynamical Mordell–Lang Conjecture}}\label{30006813}
\subsection{Definitions and mathematical statement}
For a morphism $f:X\to X$ of a quasiprojective variety over a characteristic-zero field $K$, a closed subvariety $V$, and $x\in X(K)$, define the return-time set
\[
 N_f(x,V)=\{n\in{\mathbb{Z}}_{\ge0}:f^n(x)\in V(K)\}.
\]
The conjecture says this is a finite union of sets $a+b{\mathbb{Z}}_{\ge0}$ with $a,b\ge0$. The case $b=0$ is a singleton; an empty union is allowed. Thus exceptional isolated returns and infinite periodic strings of returns are both permitted. Since $f$ is everywhere defined, all iterates exist. No dominance, smoothness, invertibility, or \"etaleness is assumed in the universal statement.
\par\smallskip\textit{Formulation and concept references:} \href{https://arxiv.org/pdf/2602.08608v2}{[S1]}.

\subsection{Short English statement}
Must the times at which an algebraic orbit visits a subvariety consist of finitely many arithmetic progressions and isolated times?
\subsection{Sources}
\begin{itemize}
\item[S1]Dynamical Mordell–Lang conjecture for split self-maps of affine curve times projective curve; arXiv2602.08608v2,16April2026.
\url{https://arxiv.org/pdf/2602.08608v2}.
\textit{Formulation source.}
\end{itemize}
\end{document}
